% !TeX root = ../../Book2.tex
\chapter{\texorpdfstring{\(L^2\)}{L²} 与 \texorpdfstring{\(L^p\)}{Lᵖ} Fourier 理论}
\label{b2:ch:04}
\BookChapterNumber{b2:ch:04}
% Axler2020Measure, Section11C, 印刷375--377页核对L2延拓的定义和证明路线。
% Hausdorff--Young在本章用有限简单函数和已证明三线引理完整推导；不前借第17章。
\begin{chapterabstract}
Fourier 变换不仅处理可积函数，还可以通过能量等式作用在整个平方可积空间。本章证明 Plancherel 延拓的良定义、满射性与反演，再用两端估计和复参数插值建立 Hausdorff--Young 不等式。空间伸缩与分离波包的例子显示，大于二的指数不能按同样方式延拓到共轭指数空间。最后以频率乘子说明酉性怎样简化线性算子，并引出广义变换的需要。
\end{chapterabstract}
\begin{learninggoals}
\begin{itemize}
\item\label{b2:item:04-goal-extension}从共同稠密类构造唯一的平方范数变换，说明定义与逼近列无关。
\item\label{b2:item:04-goal-unitary}用反演证明酉性，区分等距、满射及逐点积分三个不同问题。
\item\label{b2:item:04-goal-interpolation}理解两端估计怎样经参数函数产生中间指数，并核对插值的每个积分和极限。
\item\label{b2:item:04-goal-failure}用尺度及符号平均构造 \(p>2\) 的范数障碍，准确表述其否定范围。
\item\label{b2:item:04-goal-multiplier}在 \(L^2\) 上定义频率乘子，计算范数，并辨认推广到其他指数所需的额外信息。
\end{itemize}
\end{learninggoals}

本章接续\chapref{b2:ch:03}的反演、Gauss 核与 Young 不等式。证明 Plancherel 时，必须先在可作反演的函数上建立等式，再用单调收敛扩大范围；证明插值时，则先只处理有限测度支集的简单函数。两处都先解决定义与极限，再写空间之间的映射。

\ProseParagraphBreak

\secref{b2:sec:0401}是能量方法的核心；\secref{b2:sec:0402}独立证明本章所需的插值结论，不要求先读第三卷复分析或本卷\chapref{b2:ch:17}。\secref{b2:sec:0403}说明所得结论的适用边界。Sheldon Axler 的《Measure, Integration \& Real Analysis》\cite[Section 11C, pp. 375--377]{Axler2020Measure}可用于伴读 Plancherel 延拓；进一步的 \(L^p\) 方法可参阅 Loukas Grafakos 的《Classical Fourier Analysis》\cite{Grafakos2014Classical}。

% !TeX root = ../../Book2.tex
\section{Plancherel 定理}
\label{b2:sec:0401}

上一章把 Fourier 变换定义为绝对收敛积分，因此要求 \(f\in L^1\)。平方可积函数未必可积，但其能量仍有限。本节先在两个空间的交集上证明能量守恒，再利用完备性延拓变换。延拓所得对象是一个 \(L^2\) 等价类，不能自动解释为每个频率处的普通积分。

\subsection{\texorpdfstring{\(L^1\cap L^2\)}{L¹∩L²}}
\label{b2:subsec:040101}

先在同时具有可积变换的函数上计算。

\begin{lemma}\label{b2:lem:04-parseval-core}
若 \(f,g\in L^1\cap L^2\)，且 \(\widehat f,\widehat g\in L^1\)，则
\[
 \int f(x)\cdot\overline{g(x)}\,\mathrm{d}x
 =\int\widehat f(\xi)\cdot\overline{\widehat g(\xi)}\,\mathrm{d}\xi.
\]
\end{lemma}
\begin{proof}
反演使 \(f(x)=\int\widehat f(\xi)e^{2\pi\mathrm{i}x\cdot\xi}\,\mathrm{d}\xi\) 几乎处处。将此式乘以 \(\overline g(x)\) 再积分；重积分绝对值不超过 \(\|\widehat f\|_1\cdot\|g\|_1\)，故可交换次序。内层积分为 \(\overline{\widehat g(\xi)}\)。右侧绝对收敛还可由 \(\widehat f\in L^1\)、\(\widehat g\in L^\infty\) 直接确认。
\end{proof}

\begin{theorem}[Plancherel，交空间形式]\label{b2:thm:04-plancherel-core}
若 \(f\in L^1(\mathbb R^d)\cap L^2(\mathbb R^d)\)，则 \(\widehat f\in L^2\)，且 \(\|\widehat f\|_2=\|f\|_2\)。
\end{theorem}
\begin{proof}
令 \(f_\varepsilon=G_\varepsilon*f\)。Young 不等式保证 \(f_\varepsilon\in L^1\cap L^2\)，且 \(f_\varepsilon\to f\) 于 \(L^2\)。卷积定理给出
\[
 \widehat {f_\varepsilon}(\xi)
 =e^{-\pi\varepsilon^2|\xi|^2}\cdot\widehat f(\xi).
\]
右侧是有界函数与可积 Gauss 因子的乘积，所以属于 \(L^1\)，也属于 \(L^2\)。对 \(f_\varepsilon\) 使用前一引理，得到
\[
 \|f_\varepsilon\|_2^2
 =\int e^{-2\pi\varepsilon^2|\xi|^2}|\widehat f(\xi)|^2\,\mathrm{d}\xi.
\]
令 \(\varepsilon\downarrow0\)，左侧趋于 \(\|f\|_2^2\)；右侧由单调收敛趋于 \(\int|\widehat f|^2\)，即使预先尚不知道此积分有限也可使用。于是同时得到可积性和等式。
\end{proof}

这称为\term[plancherel-dingli]{Plancherel 定理}[Plancherel's theorem]。证明中特意先使用 Gauss 因子使变换可积；不能对一般 \(L^1\cap L^2\) 函数直接写反演，再把待证明的平方可积性用于换序。

\subsection{Plancherel 延拓}
\label{b2:subsec:040102}

\begin{theorem}\label{b2:thm:04-plancherel-extension}
存在唯一线性等距映射 \(\mathcal F_2:L^2(\mathbb R^d)\to L^2(\mathbb R^d)\)，在 \(L^1\cap L^2\) 上与可积 Fourier 变换一致。
\end{theorem}
\begin{proof}
对 \(f\in L^2\)，取 \(f_n=f\cdot\mathbf1_{\{|x|\leq n,\ |f|\leq n\}}\)。它有界、有限测度支集，故属于 \(L^1\cap L^2\)，并由控制收敛在 \(L^2\) 中趋于 \(f\)。Plancherel 等式说明 \((\widehat f_n)\) 是 Cauchy 列，定义 \(\mathcal F_2f\) 为其 \(L^2\) 极限。

若使用另一逼近列 \(g_n\)，则 \(\|\widehat f_n-\widehat g_n\|_2=\|f_n-g_n\|_2\to0\)，所以定义不依赖选择。对和与标量倍数选择相应的逼近列，便知映射线性；取范数极限得到等距性。任何连续延拓在稠密子集上相等，就在整个空间相等，故唯一。
\end{proof}
\symidx{\(\mathcal F_2\)：平方可积函数的 Fourier 变换}

当不致混淆时仍记 \(\widehat f=\mathcal F_2f\)。例如一维函数 \((1+|x|)^{-3/4}\) 平方可积却不可积；对它不能用原来的 Lebesgue 积分定义变换，因为振荡因子模为 \(1\)，不改变绝对可积性。它的变换由上述范数极限定义。

\subsection{酉性}
\label{b2:subsec:040103}

\begin{theorem}[Plancherel，酉形式]\label{b2:thm:04-unitary}
映射 \(\mathcal F_2\) 满射并保持内积，因而为\term[you-suanzi]{酉算子}[unitary operator]。若 \(Rf(x)=f(-x)\)，则
\[
 \mathcal F_2^2=R,\quad \mathcal F_2^4=I,
 \quad \mathcal F_2^{-1}=R\mathcal F_2=\mathcal F_2R.
\]
\end{theorem}
\begin{proof}
若 \(f\in C_c^\infty\)，则 \(\widehat f\in L^1\cap L^2\)，由反演有 \(\mathcal F_2^2f=Rf\)。光滑紧支集函数在 \(L^2\) 中稠密，且两边都是连续算子，故等式推广到全部 \(L^2\)。于是 \(\mathcal F_2^4=R^2=I\)，特别每个 \(g\) 都是 \(\mathcal F_2(\mathcal F_2^3g)\)，所以满射。线性等距映射保持内积：将 \(f+g\)、\(f-g\)、\(f+\mathrm{i}g\)、\(f-\mathrm{i}g\) 的范数平方展开并作差，可分别恢复内积的实部和虚部。最后由平方等式及反射的交换关系得到逆算子公式；交换关系先由可积函数换元验证，再按稠密性延拓。
\end{proof}

在本书的归一化下，Fourier 变换没有丢失或放大总能量，但一般不保持某个点的值，也不保持函数的支集。酉性属于 Hilbert 空间结构，而不是逐点大小不变。

\subsection{\texorpdfstring{\(L^2\)}{L²} 反演}
\label{b2:subsec:040104}

\begin{corollary}\label{b2:cor:04-l2-inversion}
若 \(f\in L^2\)，则
\[
 f=\lim_{R\to\infty}\int_{|\xi|\leq R}\widehat f(\xi)\cdot e^{2\pi\mathrm{i}x\cdot\xi}\,\mathrm{d}\xi
\]
在 \(L^2\) 范数中成立。Gauss 正则化也有 \(\mathcal F_2^{-1}(e^{-\pi\varepsilon^2|\xi|^2}\widehat f)\to f\) 于 \(L^2\)。
\end{corollary}
\begin{proof}
有限球上的 \(\widehat f\) 由 Cauchy--Schwarz 属于 \(L^1\)，故截断反演积分有定义。其与 \(f\) 的平方范数误差恰为 \(\int_{|\xi|>R}|\widehat f|^2\)，由尾部可积性趋零。Gauss 情形同理由频率侧的控制收敛和酉性得到。
\end{proof}

任何递增穷尽空间的可测频率区域都可替代球；所需条件只是被舍弃部分的平方质量趋零。这个结论没有保证每个点上的截断积分收敛。对于随时间演化的方程，只要频率乘子有界，就可以先在 \(L^2\) 中定义演化，再另行研究点值和正则性。



\begin{theorem}[平方可积弱导数的频率刻画]\label{b2:thm:04-weak-derivative}
设 \(f\in L^2(\mathbb R^d)\)。它具有平方可积的第 \(j\) 个弱导数，当且仅当 \(\xi_j\widehat f(\xi)\in L^2\)。此时
\[
 \mathcal F_2(\partial_jf)=2\pi\mathrm{i}\xi_j\widehat f,
 \quad\|\partial_jf\|_2=2\pi\|\xi_j\widehat f\|_2.
\]
\end{theorem}
\begin{proof}
先记下双线性积分等式 \(\int\widehat u\,v=\int u\,\widehat v\)，它对 \(u,v\in L^2\) 成立：在可积平方可积的共同稠密类上由 Fubini 证明，再按 Cauchy--Schwarz 取极限。

若 \(g=\partial_jf\in L^2\)，弱导数测试式原先只用于 \(C_c^\infty\)。对 \(\psi\in C_c^\infty\)，函数 \(\widehat\psi\) 及其导数快速下降；将 \(\chi_R\widehat\psi\) 代入，再让 \(R\to\infty\)。含截断导数的误差由 \(R^{-1}\|f\|_2\cdot\|\widehat\psi\|_2\) 控制，其余项按平方可积性取极限。因此也可将 \(\widehat\psi\) 用作测试函数，得到
\[
 \int\widehat g\,\psi
 =\int g\,\widehat\psi
 =-\int f\,\partial_j\widehat\psi
 =\int\widehat f\,(2\pi\mathrm{i}\xi_j\psi).
\]
两边的局部可积函数对全部紧支集光滑函数配对相同，故几乎处处相等。这个局部唯一性可由磨光核在 Lebesgue 点的逼近直接得出。于是必要条件成立。

反之，设 \(g=\mathcal F_2^{-1}(2\pi\mathrm{i}\xi_j\widehat f)\)。对 \(\varphi\in C_c^\infty\)，用双线性 Plancherel 等式与正号变换的求导公式，
\[
 \int g\varphi
 =\int2\pi\mathrm{i}\xi_j\widehat f\,\check\varphi
 =-\int\widehat f\,\check{\partial_j\varphi}
 =-\int f\partial_j\varphi.
\]
所以 \(g\) 是弱导数。范数式由酉性得到。
\end{proof}

对多重指标迭代，可得到整数阶 Sobolev 空间的频率描述。若沿用第一卷定义 21.1.2 中按每个多重指标计一次的范数，则
\[
 \|f\|_{H^m}^2
 =\int\left(\sum_{|\alpha|\leq m}(2\pi)^{2|\alpha|}\xi^{2\alpha}\right)|\widehat f(\xi)|^2\,\mathrm{d}\xi.
\]
括号中的多项式与 \((1+4\pi^2|\xi|^2)^m\) 可用只依赖 \(d,m\) 的正数互相控制，但高阶时一般不数值相等：后者的多项式展开带有组合系数。因此用后一权重定义的范数与既有范数等价，却不能静默地把等价写成相等。

\ProseParagraphBreak

例如对 \(f\in L^2\)，令 \(\widehat u=(1+4\pi^2|\xi|^2)^{-1}\widehat f\)。所有不超过二阶的导数乘子有界，故 \(u\in H^2\)，并在弱导数意义下解 \((1-\Delta)u=f\)。这个例子只用能量估计就取得解和两阶正则性；\chapref{b2:ch:05}将进一步把同一分母解释成基本解的变换。

% !TeX root = ../../Book2.tex
\section{Hausdorff--Young 不等式}
\label{b2:sec:0402}

变换已经有两个端点估计：可积函数变为有界函数，平方可积函数保持平方范数。中间指数的控制将通过一个连续参数把这两端连起来。这里给出 Fourier 变换所需的完整插值论证，一般算子形式留在\secref{b2:sec:1704}。

\subsection{\texorpdfstring{\(L^1\to L^\infty\)}{L¹→L∞}}
\label{b2:subsec:040201}

由积分三角不等式，\(\|\widehat f\|_\infty\leq\|f\|_1\)。该端点常数为 \(1\)，非负函数在零频率处取等。它强调总质量对每个频率的控制；若只知道 \(L^1\) 范数，就没有进一步统一的频率衰减率。

\subsection{\texorpdfstring{\(L^2\to L^2\)}{L²→L²}}
\label{b2:subsec:040202}

Plancherel 给出另一个端点 \(\|\mathcal F_2f\|_2=\|f\|_2\)。两种定义在交集上一致，才有可能对同一个算子作插值。对有限测度支集的简单函数，它既是绝对收敛积分，也是平方范数变换；下面先在这类共同对象上完成估计。

\subsection{Riesz--Thorin 思想}
\label{b2:subsec:040203}

先准备所需的复参数估计。这里的整函数只用于一个辅助引理；下面出现的函数都是有限指数和，因而可由处处收敛的幂级数直接理解。

\begin{lemma}[三线引理的有界形式]\label{b2:lem:04-three-lines}
设整函数 \(F\) 在闭带 \(0\leq\operatorname{Re}z\leq1\) 上有界，并满足 \(|F(\mathrm{i}t)|\leq1\)、\(|F(1+\mathrm{i}t)|\leq1\)。则 \(|F(\theta)|\leq1\) 对 \(0\leq\theta\leq1\) 成立。
\end{lemma}
\begin{proof}
先回顾幂级数的最大模原理。在某个内点达到模的局部最大值时，把函数乘以常数相位，使中心值为非负实数。幂级数在小圆上的平均等于中心值；而圆上实部不超过该最大值，所以连续的非负差只能处处为零。圆上函数遂为同一常数，逐频率积分又使所有正阶幂级数系数为零。函数在邻域内常值，并由重叠圆盘延拓到连通域。因此非恒定解析函数在紧矩形上的最大模出现在边界。

给定 \(\varepsilon>0\)，考虑 \(F(z)e^{\varepsilon z^2}\)。在两条竖边上，其模至多为 \(e^\varepsilon\)。若 \(|F|\leq M\)，在高度为 \(\pm R\) 的两条横边上，模至多为 \(Me^{\varepsilon(1-R^2)}\)，对充分大的 \(R\) 也不超过 \(e^\varepsilon\)。在矩形上用刚证明的最大模原理，得 \(|F(\theta)|e^{\varepsilon\theta^2}\leq e^\varepsilon\)。最后令 \(\varepsilon\downarrow0\)。
\end{proof}

这一估计是\term[sanxian-yinli]{三线引理}[three-lines lemma]的特殊形式。它把两条边界上的控制传递到中间，下面让边界对应两种不同的积分指数。

\begin{theorem}[Hausdorff--Young]\label{b2:thm:04-hausdorff-young}
设 \(1\leq p\leq2\)，\(q=p'\)，即 \(1/p+1/q=1\)。若 \(f\) 为有界且有限测度支集的简单函数，则
\[
 \|\widehat f\|_q\leq\|f\|_p.
\]
\end{theorem}
\begin{proof}
端点已证，以下设 \(1<p<2\)，并令 \(\theta=2(1-1/p)\)。先归一化 \(\|f\|_p=1\)，再取任意有限测度支集的简单函数 \(g\)，满足 \(\|g\|_p=1\)。记非零复数 \(w\) 的相位为 \(w/|w|\)，零处定义为零。设
\[
 f_z(x)=\operatorname{phase}(f(x))\,|f(x)|^{p(1-z/2)},
 \quad g_z(\xi)=\operatorname{phase}(g(\xi))\,|g(\xi)|^{p(1-z/2)},
\]
零值处仍取零，复幂使用正数 \(|f|\) 的实对数。由于函数各只有有限个非零取值，
\[
 F(z)=\int\!\int e^{-2\pi\mathrm{i}x\cdot\xi}f_z(x)\cdot g_z(\xi)\,\mathrm{d}x\,\mathrm{d}\xi
\]
是有限指数和。其指数中的系数为实数，故在闭带上有界；每个积分的绝对可积性由两个有限测度支集保证。

当 \(\operatorname{Re}z=0\) 时，\(\|f_z\|_1=\|g_z\|_1=1\)，由 \(L^1\to L^\infty\) 端点知 \(|F(z)|\leq1\)。当 \(\operatorname{Re}z=1\) 时，\(\|f_z\|_2=\|g_z\|_2=1\)，由 Plancherel 和 Cauchy--Schwarz 也有 \(|F(z)|\leq1\)。三线引理给出 \(|F(\theta)|\leq1\)；而 \(p(1-\theta/2)=1\)，所以
\[
 \left|\int\widehat f(\xi)g(\xi)\,\mathrm{d}\xi\right|\leq1.
\]
对这些 \(g\) 取上确界即得 \(\|\widehat f\|_q\leq1\)。这里即使尚未知道 \(\widehat f\in L^q\)，对偶范数判别仍合法：先在有限测度区域、有限幅值范围内取其共轭相位乘以 \(|\widehat f|^{q-1}\)，再用简单函数逼近并归一化；得到每个截断的 \(L^q\) 范数不超过 \(1\)，最后用单调收敛。去掉归一化，便得一般结论。
\end{proof}

上述证明体现\term[riesz-thorin-chazhi]{Riesz--Thorin 插值}[Riesz--Thorin interpolation]的方法：变动的是函数幅值的幂，使两个边界恰好落在已知空间。结果称为\term[hausdorff-young-budengshi]{Hausdorff--Young 不等式}[Hausdorff--Young inequality]。本章得到可用常数 \(1\)，不宣称中间指数的最佳常数也等于 \(1\)。

\subsection{\texorpdfstring{\(1\leq p\leq2\)}{1≤p≤2}}
\label{b2:subsec:040204}

\begin{corollary}\label{b2:cor:04-lp-extension}
对每个 \(1<p<2\)，Fourier 变换唯一延拓为连续线性映射 \(\mathcal F_p:L^p\to L^{p'}\)，且 \(\|\mathcal F_pf\|_{p'}\leq\|f\|_p\)。它在 \(L^1\cap L^p\) 上与可积变换一致，在 \(L^2\cap L^p\) 上与 Plancherel 变换一致。
\end{corollary}
\begin{proof}
取简单函数在 \(L^p\) 中逼近，前一定理使变换在 \(L^{p'}\) 中为 Cauchy 列。其极限不依赖逼近列，线性、连续性和唯一性与 \(L^2\) 延拓的证明相同。对交空间中的函数，可先同时截断空间与幅值，再用简单函数同时逼近两种范数。在 \(L^1\) 一端，变换一致收敛；在 \(L^2\) 一端，按平方范数收敛。有限测度区域上的依测度极限唯一，故新延拓与原变换几乎处处一致。
\end{proof}

因此在共同定义域上可以继续使用 \(\widehat f\) 这个记号。若 \(1<p<2\)，变换属于 \(L^{p'}\)，但 \(p'>2\)，不能简单地再对变换使用同一个 Hausdorff--Young 定理来反演；反演的正确表达要看它是否另外满足可积性、平方可积性，或采用后面的分布框架。



\ProseParagraphBreak

一个新定义的必要性，可以从同时不属于两个旧空间的函数看出。固定 \(1<p<2\)，选
\[
 d/2<\alpha<d/p<\beta<d,
\]
并取 \(f(x)=|x|^{-\alpha}\mathbf1_{\{|x|<1\}}+(1+|x|)^{-\beta}\mathbf1_{\{|x|\geq1\}}\)。原点附近由 \(\alpha p<d\) 保证 \(p\) 次可积，无穷远由 \(\beta p>d\) 保证，所以 \(f\in L^p\)。但原点奇性使它不平方可积，无穷远尾部使它不可积。因而它既不能整体代入\chapref{b2:ch:03}的绝对积分，也不能整体代入 Plancherel 定义；Hausdorff--Young 延拓确实增加了对象。

\ProseParagraphBreak

定义这个变换时，只需取在 \(L^p\) 中逼近 \(f\) 的简单函数列。其频率像由不等式组成 \(L^{p'}\) Cauchy 列，不必事先猜测极限的点值公式。若另外将 \(f\) 拆成可积部分与平方可积部分，两个变换的和也应与新定义相容；这种相容性将在下一章的测试配对中统一表达。

% !TeX root = ../../Book2.tex
\section{\texorpdfstring{\(L^p\)}{Lᵖ} Fourier 变换的局限}
\label{b2:sec:0403}

把指数从 \(2\) 往两侧移动，会改变变换所能取值的空间。下面用尺度和相互分离的波包解释为什么 \(p>2\) 不能照搬上一节的结论，并为分布意义下的变换留下明确的问题。

\subsection{\texorpdfstring{\(p>2\)}{p>2}}
\label{b2:subsec:040301}

先作尺度检验。若存在 \(\|\widehat f\|_q\leq C\|f\|_p\) 对全部光滑紧支集函数成立，则对 \(f(ax)\) 使用伸缩公式，得到左右尺度分别为 \(a^{-d+d/q}\) 与 \(a^{-d/p}\)。让 \(a\) 趋于零或无穷，必须有 \(1/p+1/q=1\)。这个必要条件仍不充分。

\begin{theorem}\label{b2:thm:04-p-above-two}
若 \(2<p\leq\infty\)、\(q=p'\)，则不存在常数 \(C\) 使 \(\|\widehat f\|_q\leq C\|f\|_p\) 对所有 \(f\in C_c^\infty(\mathbb R^d)\) 成立。
\end{theorem}
\begin{proof}
固定非零 \(\varphi\in C_c^\infty\)，选 \(L\) 足够大，使 \(\varphi(x-jLe_1)\) 的支集互不相交。对符号 \(\epsilon_j\in\{-1,1\}\)，令
\[
 f_N(x)=\sum_{j=1}^N\epsilon_j\varphi(x-jLe_1).
\]
于是 \(\|f_N\|_p=N^{1/p}\|\varphi\|_p\)，\(p=\infty\) 时取 \(N^{1/p}=1\)；其变换为 \(\widehat\varphi(\xi)\sum_{j=1}^N\epsilon_je^{-2\pi\mathrm{i}jL\xi_1}\)。

把全部 \(2^N\) 种符号等权平均。对任意模为 \(1\) 的复数 \(z_j\)，记 \(Z=\sum_j\epsilon_jz_j\)、\(X=|Z|^2\)。展开后，含某个符号奇数次的项平均为零，所以 \(\mathbb EX=N\)；四次矩只有指标成对出现的项保留，三个配对方式各至多贡献 \(N^2\)，故 \(\mathbb EX^2\leq3N^2\)。令 \(E=\{X\geq N/2\}\)，则
\[
 N=\mathbb EX\leq N/2+(\mathbb EX^2)^{1/2}\mathbb P(E)^{1/2},
\]
因此 \(\mathbb P(E)\geq1/12\)，进而 \(\mathbb E|Z|^q\geq12^{-1}(N/2)^{q/2}\)。

逐频率使用这个估计，乘以 \(|\widehat\varphi(\xi)|^q\) 再积分。有限平均可与积分交换，故至少有一组选定符号使
\[
 \|\widehat {f_N}\|_q\geq c_qN^{1/2}\|\widehat\varphi\|_q.
\]
光滑紧支集函数的变换快速衰减且非零，故右侧范数有限而正。若所求估计成立，则 \(N^{1/2-1/p}\) 应一致有界，与 \(p>2\) 矛盾。
\end{proof}

这个反例甚至使用光滑紧支集函数，故障碍来自范数估计本身，并非某个积分未定义。它否定的是在整个 \(L^p\) 上取值于 \(L^{p'}\) 的有界延拓；并没有否定特殊函数仍可能有很好的变换。

\subsection{点态意义}
\label{b2:subsec:040302}

在无限测度空间上，较大的 \(p\) 不意味着 \(L^1\)。例如 \((1+|x|)^{-\alpha}\) 在 \(\alpha p>d\) 时属于 \(L^p\)，但若 \(\alpha\leq d\) 就不可积。乘上单位模复指数以后，绝对可积性没有改善，所以定义式不能作为普通 Lebesgue 积分使用。

\ProseParagraphBreak

有时可以指定对称截断、Gauss 正则化或其他极限。不同极限的存在性各需证明；一般的 \(L^p\) 等价类也没有预先规定的点值。即使在 \(L^2\) 中，构造的是频率侧的等价类；一句“取每个频率的极限”并不能替代范数完备性给出的定义。

\subsection{乘子问题}
\label{b2:subsec:040303}

\begin{definition}\label{b2:def:04-multiplier}
对 \(m\in L^\infty(\mathbb R^d)\)，在 \(L^2\) 上定义
\[
 T_mf\coloneqq\mathcal F_2^{-1}(m\widehat f),
\]
称为符号 \(m\) 的\term[fourier-chengzi-suanzi]{Fourier 乘子算子}[Fourier multiplier operator]。
\end{definition}
\symidx{\(T_m\)：符号为 \(m\) 的 Fourier 乘子算子}

\begin{proposition}\label{b2:prop:04-multiplier-norm}
有 \(\|T_m\|_{L^2\to L^2}=\|m\|_\infty\)。若 \(K\in L^1\)、\(m=\widehat K\)，则 \(T_mf=K*f\) 对 \(f\in L^2\) 成立；卷积公式还定义了每个 \(L^p\)、\(1\leq p\leq\infty\) 上的有界算子，范数至多为 \(\|K\|_1\)。
\end{proposition}
\begin{proof}
酉性给出上界。对任意 \(a<\|m\|_\infty\)，集合 \(\{|m|>a\}\) 有正测度，可取其中一个有限正测度子集 \(E\)，令 \(\widehat f=\mathbf1_E\)。于是 \(\|T_mf\|_2>a\|f\|_2\)，令 \(a\uparrow\|m\|_\infty\) 得到等号。卷积表示先对 \(f\in L^1\cap L^2\) 由卷积定理成立，再以 Young 估计和稠密性推广到 \(L^2\)。其余指数直接由 Young 不等式给出。
\end{proof}

对 \(p\ne2\)，仅有 \(m\) 有界通常不足以取得同样的乘子结论；符号的平滑程度、奇性与维数会参与估计。本章不预先断言某个粗糙符号在哪些指数上有界，相关问题将在\chapref{b2:ch:22}系统研究。


\ProseParagraphBreak

平滑频率截断说明为何符号的导数有用。取 \(m\in C_c^\infty\)，令 \(K=\check m\)。由\chapref{b2:ch:03}的微分衰减估计，\(K\in L^1\)，所以 \(T_m\) 在每个 \(L^p\) 上都由 Young 不等式控制。对 \(m_R(\xi)=m(\xi/R)\)，相应核为 \(K_R(x)=R^dK(Rx)\)，且 \(\|K_R\|_1=\|K\|_1\)。于是这一族算子对全部尺度具有同一个范数界。

\ProseParagraphBreak

相反，一维的硬截断符号 \(m=\mathbf1_{[-R,R]}\) 对应
\[
 K(x)=\dfrac{\sin(2\pi Rx)}{\pi x},
\]
该核不属于 \(L^1\)。对 \(f\in C_c^\infty\)，仍可用有限频率积分和 Fubini 验证 \(T_mf=K*f\)，但不能再用 \(\|K\|_1\) 给出有限的 Young 上界。\(L^2\) 上的范数仍为 \(1\)，因为符号是有界的正交投影。这个例子把两种方法的职责分开：酉性直接控制平方能量，绝对可积核则提供全部指数上的一种充分条件。

\subsection{分布意义的 Fourier 变换的动机}
\label{b2:subsec:040304}

若不能要求变换是一个普通函数，还可以要求它对足够好的测试函数产生确定的数值。对 \(f\in L^p\) 和光滑紧支集的 \(\varphi\)，函数 \(\widehat\varphi\) 快速衰减，因而 \(\int f\widehat\varphi\) 总有意义。若 \(f\) 原来可积，Fubini 又给出
\[
 \int\widehat f(\xi)\varphi(\xi)\,\mathrm{d}\xi
 =\int f(x)\widehat\varphi(x)\,\mathrm{d}x.
\]
右侧因此提供了延拓变换的线索，且对每个 \(1\leq p\leq\infty\) 都适用。下一章将选取在 Fourier 变换下保持不变的 Schwartz 测试空间，使这个配对成为完整定义。


\begin{chaptersummary}
Plancherel 等式把变换从绝对收敛积分延拓为能量空间上的酉算子。反演和四次变换等于恒等的关系保证满射性；频率截断的误差则由被舍弃的平方质量精确给出。Hausdorff--Young 把两个端点连接起来，但其共轭指数方向只覆盖 \(1\leq p\leq2\)。

\ProseParagraphBreak

大于二时，空间中相互分离的函数可在频率中产生更大的叠加，破坏所需范数界。应当区分某个好函数的变换存在、变换在整个空间上连续，以及逐点反演这三个层次。\(L^2\) 乘子由符号本性上确界直接控制；一般指数的乘子问题，以及所有 \(L^p\) 函数的统一广义变换，则需要更大的理论框架。
\end{chaptersummary}

\BookExercises

本组题目以 Plancherel 为计算工具，分别考察尺度、频带、乘子、算子收敛及空间与频率的同时集中。第五题后的论证可以作为阅读下一章时的具体模型。

% 自拟：Gauss范数检查非最优Hausdorff--Young常数及尺度不变性。
\begin{exercise}\label{b2:ex:04-gaussian-ratio}
设 \(1<p<2\)、\(q=p'\)，\(G_a(x)=e^{-\pi a|x|^2}\)，其中 \(a>0\)。计算 \(\|G_a\|_p\)、\(\|\widehat G_a\|_q\) 和二者的比值，说明为什么只改变 \(a\) 不能改善这个比值。该计算是否已经证明最佳 Hausdorff--Young 常数？
\end{exercise}
\begin{solution}
尺度公式给出
\[
 \|G_a\|_p=(ap)^{-d/(2p)},\quad
 \widehat G_a(\xi)=a^{-d/2}e^{-\pi|\xi|^2/a},\quad
 \|\widehat G_a\|_q=a^{-d/2}(q/a)^{-d/(2q)}.
\]
因 \(1/p+1/q=1\)，比值为 \(p^{d/(2p)}q^{-d/(2q)}\)，与 \(a\) 无关。这个数是任何普遍常数的下界；只计算一族函数，不能证明所有函数都满足取此常数的上界，因而不能单凭此计算宣称最佳性。
\end{solution}

% 自拟：紧频带产生光滑代表和定量导数界。
\begin{exercise}\label{b2:ex:04-bandlimited}
设 \(f\in L^2(\mathbb R^d)\)，且 \(\widehat f\) 几乎处处支撑在 \(B(0,R)\) 内。证明 \(f\) 有光滑代表，并对每个多重指标 \(\alpha\) 有
\[
 \|\partial^\alpha f\|_\infty
 \leq(2\pi R)^{|\alpha|}|B(0,R)|^{1/2}\|f\|_2.
\]
\end{exercise}
\begin{solution}
有限频带上的 \(\widehat f\) 可积，乘任意多项式后仍可积。正号积分因此给出一个可任意次逐项求导的连续函数，且其 \(L^2\) 等价类由反演等于 \(f\)。求导后用 Cauchy--Schwarz：
\[
 |\partial^\alpha f(x)|
 \leq(2\pi R)^{|\alpha|}\int_{B(0,R)}|\widehat f|
 \leq(2\pi R)^{|\alpha|}|B(0,R)|^{1/2}\|\widehat f\|_2.
\]
最后用 Plancherel。这里的 \(R\) 有限是保证全部频率矩可积的关键。
\end{solution}

% 自拟：L2反演确实超出绝对可积函数域。
\begin{exercise}\label{b2:ex:04-sinc}
在一维令 \(g(\xi)=\mathbf1_{[-1/2,1/2]}(\xi)\)，\(f=\mathcal F_2^{-1}g\)。求 \(f\) 的连续代表，证明 \(f\in L^2\setminus L^1\)，并计算 \(\|f\|_2\)。
\end{exercise}
\begin{solution}
有限频带积分给出 \(f(x)=\sin(\pi x)/(\pi x)\)，原点补为 \(1\)。Plancherel 给出 \(\|f\|_2=\|g\|_2=1\)。在每个 \([n+1/6,n+5/6]\) 上，\(|\sin(\pi x)|\geq1/2\)，因此其绝对积分至少为常数乘 \(\sum_{n\geq1}(n+1)^{-1}\)，故不属于 \(L^1\)。这不妨碍其 \(L^2\) 变换等于区间示性函数。
\end{solution}

% 自拟：平移群强连续却不按算子范数连续。
\begin{exercise}\label{b2:ex:04-translation-operator}
证明对每个 \(f\in L^2\)，\(\|\tau_hf-f\|_2\to0\) 当 \(h\to0\)；但对每个 \(h\ne0\)，\(\|\tau_h-I\|_{L^2\to L^2}=2\)。
\end{exercise}
\begin{solution}
平移的变换公式按稠密性推广到 \(L^2\)。于是
\[
 \|\tau_hf-f\|_2^2
 =\int|e^{-2\pi\mathrm{i}h\cdot\xi}-1|^2|\widehat f(\xi)|^2\,\mathrm{d}\xi\longrightarrow0
\]
由控制收敛成立。乘子范数则等于 \(e^{-2\pi\mathrm{i}h\cdot\xi}-1\) 的本性上确界；\(h\ne0\) 时，可在正测度频率带上使相位任意接近 \(-1\)，故该本性上确界为 \(2\)。前一结论固定函数，后一结论允许随 \(h\) 改变测试函数。
\end{solution}

% 自拟：频率投影和正性、伴随的分类。
\begin{exercise}\label{b2:ex:04-multiplier-classification}
设 \(m\in L^\infty\)。求 \(T_m^*\)，并分别刻画 \(T_m\) 自伴、非负、酉和正交投影的充要条件。这里非负指 \(\langle T_mf,f\rangle\geq0\) 对全部 \(f\in L^2\) 成立。
\end{exercise}
\begin{solution}
在频率侧计算内积，得 \(T_m^*=T_{\overline m}\)。所以自伴当且仅当 \(m\) 几乎处处为实数。非负当且仅当 \(m\geq0\) 几乎处处：充分性由积分给出；若实部在某个正测度集合上为负，或虚部不为零，可取该集合有限测度部分的示性函数作频率测试，违反非负实值条件。酉性等价于 \(|m|=1\) 几乎处处，此时逆符号为 \(\overline m\)。正交投影等价于自伴且幂等，故 \(m\) 几乎处处只能取 \(0,1\)，即为某个可测集的示性函数。
\end{solution}

% 自拟：傅里叶算子的四个代数谱分量；无须预借Hermite函数。
\begin{optionalexercise}\label{b2:ex:04-four-sectors}
对 \(\lambda\in\{1,-1,\mathrm{i},-\mathrm{i}\}\)，定义
\[
 Q_\lambda=\dfrac14\sum_{j=0}^3\lambda^{-j}\mathcal F_2^j.
\]
证明这些算子是两两正交的正交投影，其和为 \(I\)，且其值域等于 \(\{f:\mathcal F_2f=\lambda f\}\)。说明 \(\lambda=\pm1\) 的值域只含偶函数，\(\lambda=\pm\mathrm{i}\) 的值域只含奇函数。
\end{optionalexercise}
\begin{solution}
利用 \(\mathcal F_2^4=I\) 移动求和指标，可得 \(\mathcal F_2Q_\lambda=\lambda Q_\lambda\)。若 \(\mathcal F_2f=\mu f\)，则 \(Q_\lambda f=\frac14\sum_{j=0}^3(\mu/\lambda)^jf\)，当 \(\mu=\lambda\) 时等于 \(f\)，否则为零。并且四个 \(Q_\lambda\) 的和为 \(I\)，因为四次单位根的非零次幂和为零。这已给出幂等性、相互乘积为零和值域描述。酉性使不同单位模特征值的特征向量正交：\(\langle f,g\rangle=\lambda\overline\mu\langle f,g\rangle\)。故这些投影确为正交投影。最后 \(Rf=\mathcal F_2^2f=\lambda^2f\)，得到奇偶性。
\end{solution}

% 自拟：有限空间与频率支集的基本体积障碍。
\begin{exercise}\label{b2:ex:04-support-volume}
设非零 \(f\in L^2\) 分别在空间与频率侧几乎处处支撑于有限测度集 \(A,B\)。证明 \(|A|\cdot|B|\geq1\)。
\end{exercise}
\begin{solution}
空间支集有限使 \(f\in L^1\)，且 \(\|f\|_1\leq|A|^{1/2}\|f\|_2\)。于是
\[
 \|f\|_2^2=\|\widehat f\|_2^2
 \leq|B|\cdot\|\widehat f\|_\infty^2
 \leq|B|\cdot\|f\|_1^2
 \leq|A|\cdot|B|\cdot\|f\|_2^2.
\]
因为 \(f\ne0\)，除去正的平方范数即得所求。该结论只是一个必要条件，没有断言满足体积条件的任意两个集合都能成为精确双侧支集。
\end{solution}

% 自拟：把精确支集障碍改成有误差的集中估计。
\begin{optionalexercise}\label{b2:ex:04-approximate-support}
设 \(\|f\|_2=1\)，\(A,B\) 有限测度，并有 \(\|f\cdot\mathbf1_{A^c}\|_2\leq\varepsilon\)、\(\|\widehat f\cdot\mathbf1_{B^c}\|_2\leq\delta\)，其中 \(\varepsilon+\delta<1\)。证明 \(|A|^{1/2}\cdot|B|^{1/2}\geq1-\varepsilon-\delta\)。
\end{optionalexercise}
\begin{solution}
令 \(P_Af=\mathbf1_A\cdot f\)、\(Q_Bf=\mathcal F_2^{-1}(\mathbf1_B\cdot\widehat f)\)。两者均为正交投影，故
\[
 \|f-Q_BP_Af\|_2
 \leq\|f-Q_Bf\|_2+\|Q_B(f-P_Af)\|_2\leq\delta+\varepsilon.
\]
另一方面，\(P_Af\in L^1\)，所以与前题相同的估计给出
\[
 \|Q_BP_Af\|_2
 =\|\mathbf1_B\cdot\widehat{P_Af}\|_2
 \leq|B|^{1/2}\|P_Af\|_1
 \leq|A|^{1/2}\cdot|B|^{1/2}.
\]
由三角不等式，\(1\leq\varepsilon+\delta+\|Q_BP_Af\|_2\)，得到所求。
\end{solution}
